\subsection{Leaky-Phase Analysis Under the Locked Invariant Class}
\label{sec:leaky_phase_analysis}

We analyze the Paper-1 invariant class
\[
\frac{dm}{d\tau}=\beta-\lambda m,\qquad
\Delta P = k\,[m(\tau_T)-m(\tau_0)].
\]
The decay parameter $\lambda$ controls the relaxation rate of $m(\tau)$ but cannot create matched-endpoint splitting under constant $\tau$-intrinsic drive; it only attenuates splitting already induced by nontrivial $\tau$-intrinsic inputs.
For fixed endpoints $(\tau_0,\tau_T)$ and fixed initial condition, $\Delta P$ is endpoint-determined; matched-endpoint path pairs therefore do not split under this class.

\paragraph{Numerical phase behavior (\texttt{paper\_final\_20260302}).}
From \texttt{tables/leaky\_memory\_grid.csv} (rows with $\beta=1$, $\mathrm{dt\_hint}=0.5$), terminal memory decreases monotonically with $\lambda$:
\[
m_{\text{end}}(10^{-4})=4256.584,\;
m_{\text{end}}(10^{-3})=996.1074,\;
m_{\text{end}}(10^{-2})=100,\;
m_{\text{end}}(10^{-1})=10,\;
m_{\text{end}}(1)=1.
\]
The matched-path endpoint difference remains numerically zero across the tested grid:
\[
\Delta m_{\text{end}}=0,\qquad \Delta P_{\text{path}}=0.
\]

\paragraph{Interpretation.}
The leaky parameter $\lambda$ controls memory amplitude and horizon, not gauge violation. In this locked invariant class:
\begin{itemize}
\item $\lambda$ modulates endpoint state magnitude $m(\tau_T)$.
\item matched-endpoint path comparisons remain collapsed ($\Delta P_{\text{path}}=0$).
\item gauge separation must therefore be established against coordinate-dependent observers (Section~\ref{sec:model_class_comparison}).
\end{itemize}
